\chapter{Cyst}

\section*{Definition}
A cyst is a closed capsule or sac-like structure filled with liquid, semisolid, or gaseous material. Cysts can develop in any tissue and vary widely in size. Most cysts are benign, though some require evaluation to exclude malignancy.

\section*{Pathophysiology}
Cysts form when fluid-producing cells become enclosed within a membranous sac, or when a duct becomes blocked and secretions accumulate behind the obstruction. Common types include epidermal (sebaceous) cysts from obstructed hair follicles, ganglion cysts from synovial fluid herniation, ovarian cysts, Baker's cysts, and breast cysts.

\section*{Causes}
\begin{itemize}
\item Blocked ducts or glands (sebaceous gland obstruction)
\item Trauma (inclusion cysts)
\item Congenital remnants (branchial cleft cyst, thyroglossal duct cyst)
\item Infection (pilonidal cyst)
\item Hormonal factors (ovarian follicles, breast cysts)
\item Joint or tendon pathology (ganglion cyst, Baker's cyst)
\item Tumors (cystic degeneration of neoplasms)
\item Inflammatory conditions (ovarian endometrioma)
\end{itemize}

\section*{When to See a Doctor}
\begin{itemize}
\item Palpable lump under the skin or in deeper tissue
\item Smooth, round, mobile mass
\item Slow or no growth over time
\item Pain, tenderness, or redness if inflamed or infected
\item Rapid growth or change in character
\item Cyst in unusual location
\item Cyst causing functional impairment (pain, pressure on nerves)
\end{itemize}

\section*{Related Symptoms}
\begin{itemize}
\item Abscess (infected, pus-filled cavity)
\item Lipoma (benign fatty tumor)
\item Tumor (solid neoplasm versus cystic structure)
\item Boil (furuncle, infected hair follicle)
\item Lymphedema (swelling from lymphatic obstruction)
\item Hematoma (blood collection)
\end{itemize}

\section*{Self-Care/Home Management}
\begin{itemize}
\item Do NOT attempt to squeeze, puncture, or drain a cyst at home
\item Apply warm compresses if the cyst is inflamed
\item Keep the area clean to prevent infection
\item Monitor for changes in size, color, or pain level
\item Avoid friction or pressure over the cyst
\item Seek medical evaluation if the cyst becomes painful, red, or grows rapidly
\end{itemize}

\section*{How Your Doctor Will Evaluate You}
\begin{itemize}
\item Clinical examination (palpation, transillumination)
\item Ultrasound to determine cystic vs. solid nature
\item CT or MRI for deeper cysts
\item Fine-needle aspiration for cytology if malignancy suspected
\item Excisional biopsy for definitive diagnosis
\end{itemize}

\section*{Treatment Options}
\begin{itemize}
\item Observation if asymptomatic and benign-appearing
\item Incision and drainage for infected cysts
\item Excision (complete removal of cyst and lining) to prevent recurrence
\item Aspiration for simple cysts (higher recurrence rate)
\item Sclerotherapy for some types (ganglion, ovarian)
\item Antibiotics for secondary infection
\item Treat underlying condition
\end{itemize}
